% Source: https://artofproblemsolving.com/community/c3h2388566_symmetric_equations
% Stable identity: aops-topic-2388566-q01
\textbf{Symmetry.} Exchanging $p$ and $q$ turns the equation into $q^2+p^2=qpn+1$, which is the same equation. So $(p,q,n)$ is a solution exactly when $(q,p,n)$ is. Hence it suffices to find the solutions with $p\le q$; the others are obtained by swapping $p$ and $q$.

\textbf{Case $p=q$.} Then $2p^2-1=np^2$, so $p^2$ divides $1$, which is impossible. So $p<q$.

\textbf{Case $p<q$.} Reducing the equation modulo $q$ gives $p^2\equiv1\pmod q$, i.e. $q\mid(p-1)(p+1)$. As $q$ is prime, $q\mid p-1$ or $q\mid p+1$. But $1\le p-1<q$, so $q\nmid p-1$. Hence $q\mid p+1$, and since $q>p$, i.e. $q\ge p+1$, this forces $q=p+1$. The only two primes that differ by $1$ are $2$ and $3$, so $p=2$, $q=3$. Then $4+9=6n+1$ gives $n=2$.

So the solutions are $\boxed{(p,q,n)=(2,3,2)\ \text{and}\ (3,2,2)}$.
